\section{Experimental setup and validation}
All test conditions are unseen by the fitted models. Protocols: models fitted on 25\,°C and tested on 0\,°C, 45\,°C and 1\,A; two-temperature training with the third temperature held out; all three low-current tests used to predict the 1\,A test (unseen rate). The first 10\,min of each test are excluded for every model so that all are scored on identical samples. Hyper-parameters of tree models are chosen by blocked cross-validation (5 contiguous blocks per training condition, gap equal to the window length) so overlapping samples never straddle the train/validation boundary; recurrent networks use early stopping on a contiguous validation block. Seeds are fixed (42). Noise robustness uses \emph{simulated} perturbations added to the real test signals (Table~\ref{tab:noisedef}). After observing that LSTM/GRU collapse under noise, a variant with a floor on the input standard deviation (\texttt{*\_sf}) was added post hoc; the original variants are kept in every table.
\begin{table}[H]\centering\small
\begin{tabular}{lrrr}\toprule Level & $\sigma_V$ (mV) & $\sigma_I$ (mA) & $I$ bias (mA)\\\midrule low & 5 & 10 & 0\\ medium & 10 & 25 & 0\\ high & 20 & 50 & 0\\ I bias & 0 & 0 & 20\\\bottomrule\end{tabular}
\caption{Simulated perturbation levels.}\label{tab:noisedef}\end{table}
Blocked-CV errors of the tree models (about 7--8 SOC points) are much larger than test errors because each held-out block requires extrapolation along the SOC axis; they serve model selection only.
